\documentclass[11pt]{article}

\newcommand{\specPdfTitle}{An Algebraic Specification of fsb}
\newcommand{\specHeader}{Algebraic Specification}
% Shared by every specification in this folder. A document defines
% \specPdfTitle and \specHeader before inputting this file.
\usepackage[T1]{fontenc}
\usepackage[utf8]{inputenc}
\usepackage{lmodern}
\usepackage[margin=1in]{geometry}
\usepackage{microtype}
\usepackage{parskip}
\usepackage{amsmath}
\usepackage{amssymb}
\usepackage{amsthm}
\usepackage{mathtools}
\usepackage{booktabs}
\usepackage{longtable}
\usepackage{array}
\usepackage{enumitem}
\usepackage{needspace}
\usepackage{xcolor}
\usepackage{listings}
\usepackage{fancyhdr}
\usepackage{hyperref}
\usepackage{bookmark}
\usepackage{xurl}
\usepackage{tikz}
\usetikzlibrary{arrows.meta,positioning,fit}

\definecolor{codeblue}{HTML}{1F4E79}
\definecolor{codegreen}{HTML}{3A6B35}
\definecolor{codegray}{HTML}{555555}
\definecolor{codebackground}{HTML}{F5F7F9}
\definecolor{rulegray}{HTML}{D7DCE2}

\hypersetup{
  colorlinks=true,
  linkcolor=codeblue,
  urlcolor=codeblue,
  citecolor=codeblue,
  pdftitle={\specPdfTitle},
  pdfsubject={fsb, a read-only local file browser}
}

\lstdefinestyle{program}{
  basicstyle=\small\ttfamily,
  backgroundcolor=\color{codebackground},
  frame=single,
  rulecolor=\color{rulegray},
  framerule=0.4pt,
  xleftmargin=0.5em,
  xrightmargin=0.5em,
  aboveskip=0.8em,
  belowskip=0.8em,
  breaklines=true,
  columns=fullflexible,
  keepspaces=true,
  showstringspaces=false,
  keywordstyle=\color{codeblue}\bfseries,
  commentstyle=\color{codegreen},
  stringstyle=\color{codegray}
}

\lstset{style=program}

\newtheorem{law}{Law}[section]
\newtheorem{proposition}[law]{Proposition}
\theoremstyle{definition}
\newtheorem{definition}[law]{Definition}
\newtheorem{observation}[law]{Observation}
\newtheorem{boundary}[law]{Boundary}

\setlist{itemsep=0.25em, topsep=0.4em}
\setcounter{tocdepth}{2}
\setcounter{secnumdepth}{3}
\raggedbottom
\sloppy
\emergencystretch=3em

\pagestyle{fancy}
\fancyhf{}
\fancyhead[L]{\texttt{fsb}}
\fancyhead[R]{\specHeader}
\fancyfoot[C]{\thepage}
\setlength{\headheight}{14pt}


% Macros shared by every specification in this folder.
\newcommand{\code}[1]{\texttt{#1}}
\newcommand{\Tests}[1]{\par\smallskip\noindent\begingroup\raggedright\small\textit{Checked by:} \path{#1}\par\endgroup}

% A numbered requirement of the functional specification. The identifier is a
% label, so \rid{ID} is a link and a reference to an identifier that does not
% exist is reported by LaTeX as an undefined reference.
\newenvironment{req}[1]{\par\medskip\noindent\phantomsection\label{req:#1}\textbf{#1}\hspace{0.7em}\ignorespaces}{\par}
\newcommand{\rid}[1]{\hyperref[req:#1]{\textsf{#1}}}
% A requirement of the functional specification named from another document.
\newcommand{\fq}[1]{\textsf{#1}}
\newcommand{\shall}{\textit{shall}}

% Written by tools/stamp-version.sh; do not edit by hand.
\newcommand{\fsbversion}{1.0.11}
\newcommand{\fsbcommit}{a6d1043}
\newcommand{\fsbcodedate}{October 9, 2026}

\newcommand{\Bool}{\ensuremath{\mathsf{Bool}}}
\newcommand{\Name}{\ensuremath{\mathsf{Name}}}
\newcommand{\Path}{\ensuremath{\mathsf{Path}}}
\newcommand{\Rule}{\ensuremath{\mathsf{Rule}}}
\newcommand{\RuleSet}{\ensuremath{\mathsf{RuleSet}}}
\newcommand{\Entry}{\ensuremath{\mathsf{Entry}}}
\newcommand{\Row}{\ensuremath{\mathsf{Row}}}
\newcommand{\Type}{\ensuremath{\mathsf{Type}}}
\newcommand{\Dir}{\mathsf{dir}}
\newcommand{\File}{\mathsf{file}}

\newcommand{\Fold}{\mathsf{fold}}
\newcommand{\fold}[1]{\Fold(#1)}
\newcommand{\Hit}{\mathsf{hit}}
\newcommand{\Excl}{\mathsf{excl}}
\newcommand{\Deny}{\mathsf{Deny}}
\newcommand{\Hide}{\mathsf{Hide}}
\newcommand{\Roots}{\mathsf{Roots}}
\newcommand{\Real}{\mathsf{real}}
\newcommand{\Clean}{\mathsf{clean}}
\newcommand{\Canon}{\mathsf{canon}}
\newcommand{\InRoots}{\mathsf{inRoots}}
\newcommand{\Ok}{\mathsf{Ok}}
\newcommand{\Elig}{\mathsf{Elig}}
\newcommand{\Open}{\mathsf{Open}}
\newcommand{\List}{\mathsf{List}}
\newcommand{\Search}{\mathsf{Search}}
\newcommand{\Visible}{\mathsf{Visible}}
\newcommand{\Match}{\mathsf{Match}}
\newcommand{\MatchBelow}{\mathsf{MatchBelow}}
\newcommand{\NotFound}{\mathsf{NotFound}}
\newcommand{\Sniff}{\mathsf{sniff}}
\newcommand{\Anc}{\mathsf{anc}}
\newcommand{\Comps}{\mathsf{comps}}
\newcommand{\Parent}{\mathsf{parent}}

\newcommand{\Nfc}{\mathsf{nfc}}
\newcommand{\Sort}{\mathsf{sort}}
\newcommand{\Cmp}{\mathsf{cmp}}
\newcommand{\Filter}{\mathsf{filter}}
\newcommand{\ToHash}{\mathsf{toHash}}
\newcommand{\FromHash}{\mathsf{fromHash}}
\newcommand{\Goto}{\mathsf{goto}}
\newcommand{\Link}{\mathsf{resolveLink}}

\newcommand{\Seq}[1]{\langle #1\rangle}
\newcommand{\concat}{\mathbin{+\!+}}
\newcommand{\sqleq}{\preceq}
\newcommand{\Reach}{\rightsquigarrow}


\title{An Algebraic Specification of \texttt{fsb}}
\author{}
\date{Version \fsbversion \\ \fsbcodedate \\ \small Describes the code at commit \texttt{\fsbcommit}}

\begin{document}

\maketitle
\begin{abstract}
This document specifies the security-relevant core of \texttt{fsb}, a local,
read-only web view of a filesystem, as a small set of algebraic laws: the rule
matcher that implements the two exclusion tiers (hide and deny), the access
decision that every endpoint must agree with, the content-typed endpoints
(images, PDFs, archives, Quick Look pictures), the pure frontend functions (ordering, filtering,
navigation, link resolution), and the preview state machine.  Each law is tied
to an executable check.  Deriving the laws, and then auditing every place that compares names,
exposed sixteen defects, eight of them bypasses or gaps in the access rules, and a
later independent review of the Quick Look endpoint two more; they are recorded in
Section~\ref{sec:findings}.
\end{abstract}

\tableofcontents

\section{Purpose and Method}

\subsection{Why an algebra}

\texttt{fsb}'s central promise is negative: a path that the user has denied is
never listed, searched, previewed, downloaded, or even confirmed to exist.  That
promise is a statement about \emph{every} way of reaching a file, so it is
threatened whenever two code paths answer the same question differently.  The
defects found and fixed in the project so far had exactly that shape: a listing
through a symlinked directory that leaked denied names, a trailing newline that
defeated a rule, an unreadable file answering 403 instead of 404, and a macOS
firmlink alias that bypassed a rule.

An algebraic specification forces the question to be written down once, and
then every operation to be compared with it.  This document does that for four
parts of the system:

\begin{enumerate}
  \item the rule matcher (Section~\ref{sec:rules});
  \item the access decision shared by all endpoints (Section~\ref{sec:access});
  \item the content-typed endpoints and the pure frontend functions
        (Sections~\ref{sec:content} and~\ref{sec:frontend}); and
  \item the preview state machine (Section~\ref{sec:state}).
\end{enumerate}

\subsection{Method}

Each law is stated in the form
\emph{for all inputs of a stated kind, a stated equation or implication holds},
then given an executable check: a randomized property test with a fixed seed, a
test against a fixture that assembles the awkward cases together, or an
end-to-end browser test.  Laws that failed when first checked are recorded in
Section~\ref{sec:findings} together with the fix.  Laws that hold only under a
stated assumption say so, and the assumptions that cannot be discharged inside
the program are collected as \emph{boundaries}.

\subsection{What is not covered}

The document does not model layout, streaming, timing, or the browser's own
behavior (except where the program relies on it).  It does not replace an
independent review; it is intended to give a reviewer a precise statement of
what the program claims.

\section{Sorts and Notation}

\subsection{Sorts}

\begin{description}
  \item[$\Name$] A path component: a finite sequence of bytes containing neither
        \code{/} nor NUL.  On macOS it is always valid UTF-8, and is read as the
        sequence of Unicode characters it encodes; on Linux it may be any such bytes
        (Law~\ref{law:wire} says how the client is given one that is not UTF-8).
  \item[$\Path$] An absolute path, written as the sequence of its components
        $p=\Seq{n_1,\dots,n_k}$; the root is $\Seq{}$.  $\Comps(p)$ is that
        sequence, $\Parent(p)=\Seq{n_1,\dots,n_{k-1}}$, and $p\cdot n$ appends a
        component.
  \item[$\Type$] $\{\Dir,\File\}$, the type of the final component.
  \item[$\Rule$] A parsed pattern in the gitignore dialect, possibly negated and
        possibly restricted to directories.
  \item[$\RuleSet$] A finite sequence of rules.  Two sets are used: $\Deny$
        (hard block, no negation) and $\Hide$ (cosmetic, negation allowed).
  \item[$\Entry$] A directory entry as shown in a listing.
  \item[$\Row$] An entry as the frontend sorts it: name, type, size, and
        modification time.
  \item[$\Conn$] A TCP connection to \texttt{fsb}, given by the server's and the
        client's address and port.
  \item[$\Uid$] A user, as the system's numeric user ID.
\end{description}

\subsection{Environment}

Let $\Roots$ be the finite set of served roots, each already resolved to its
real location.  The filesystem is treated as a partial function from paths to
files and directories with symbolic links; $\Real(p)$ is the canonical real
location of $p$ after all symbolic links and aliases (such as the macOS
firmlink prefix \code{/System/Volumes/Data}) are resolved, defined when $p$
exists.  $\Canon$ reduces aliases of the same location to one spelling.

For the session: $\tau$ is the launch token of this start, $u\in\Uid$ the user
\texttt{fsb} runs as, and $\Used\in\Bool$ whether the token has been exchanged
(initially false).  $\Owner:\Conn\rightharpoonup\Uid$ is the user owning the
client end of a connection.  On Linux it is defined when the connection table
lists that end with a user \texttt{fsb}'s user namespace can name; on macOS it is
never consulted.

\subsection{Operations}

The operations analyzed are, on the server:
\[
  \Open,\ \List,\ \Search,\ \text{and the endpoints built on them,}
\]
namely file download, head, metadata, image and PDF preview, and archive
listing; the launch exchange $\Exchange$; and on the client the pure functions $\Cmp$ and $\Sort$ (row order),
$\Filter$, $\ToHash$ and $\FromHash$, $\Goto$, and $\Link$.

\subsection{Conventions}

A relation $x\approx y$ between names means \emph{the filesystem treats $x$ and
$y$ as the same name}; it is defined in Section~\ref{sec:names}.  A statement
``$\Open(p)$ fails as $\NotFound$'' means that the operation fails with an error
that the HTTP layer answers with an identical 404.

\section{Names and Folding}
\label{sec:names}

\subsection{The name equivalence}

The rules are matched against names that the operating system will treat as
identical.  On a case-insensitive APFS volume two names denote the same file
exactly when they are equal after Unicode normalization and \emph{full} case
folding.  Full folding differs from simple folding: it expands some characters
into several.

\begin{definition}[Fold]
\label{def:fold}
$\fold{s}=\Nfc\bigl(\mathsf{cf}(\Nfc(s))\bigr)$, where $\mathsf{cf}$ is Unicode
full case folding.  For names, $x\approx y \iff \fold{x}=\fold{y}$.  The
extension to paths applies $\Fold$ to the whole string, which is equivalent to
folding component by component because folding never introduces or removes a
separator.
\end{definition}

\begin{law}[Idempotence]
\label{law:fold-idempotent}
$\fold{\fold{s}}=\fold{s}$.
\end{law}

\begin{observation}[Measured aliases]
\label{obs:aliases}
The set of single characters that APFS folds onto ASCII names was measured by
creating every ASCII name of length one and two, and then probing every
character in the Unicode range \code{U+0080}--\code{U+2FFFF}.  On the tested
system the aliases are exactly the following; each row names the ASCII
spelling that the character opens.
\begin{center}
\begin{tabular}{@{}ll@{}}
\toprule
character & opens \\
\midrule
\code{\ss} (sharp s), capital sharp s (\code{U+1E9E}) & \code{ss} \\
long s (\code{U+017F}) & \code{s} \\
Kelvin sign (\code{U+212A}) & \code{k} \\
ligatures \code{U+FB00}--\code{U+FB06} & \code{ff}, \code{fi}, \code{fl}, \code{ffi}, \code{ffl}, \code{st} \\
\bottomrule
\end{tabular}
\end{center}
This is full case folding, not simple folding: under simple folding
\code{\ss} would remain a distinct character.
\end{observation}

\begin{law}[Fold agrees on the measured aliases]
\label{law:fold-aliases}
For every pair in Observation~\ref{obs:aliases},
$\fold{\text{character}}=\fold{\text{ASCII spelling}}$.  The server's $\Fold$
(Go, \code{x/text/cases}) and the client's $\Fold$ (JavaScript,
lower-, upper-, then lower-casing) both satisfy the law.
\Tests{rules/fold_test.go::TestFoldMatchesTheAliasesAPFSTreatsAsEqual; web/src/lib/laws.test.ts (filter agreement)}
\end{law}

\begin{boundary}[Folding is a property of the volume]
\label{bnd:volume}
The alias set is a fact about the file system, not about \texttt{fsb}.  A
case-sensitive volume distinguishes more names and a future OS release could
change the table.  The specification therefore treats $\approx$ as the
\emph{coarsest} equivalence the program must respect: comparing after folding
can only over-match, and over-matching a deny rule fails closed.
\end{boundary}

\section{The Rule Algebra}
\label{sec:rules}

\subsection{Definitions}

A rule $r$ has a compiled test $\Hit(r,q,t)$: whether $r$ selects the path $q$
whose final component has type $t$.  It is evaluated on folded text
(Definition~\ref{def:fold}) with a pattern that was folded the same way.

\begin{definition}[Exclusion]
\label{def:excl}
For a rule set $\mathcal R$ and a path $p$ of type $t$, let
$\Anc(p)=\{q\mid q\text{ is a prefix of }p\}$, where each proper prefix has type
$\Dir$ and $p$ has type $t$.  The last rule of $\mathcal R$ that hits $q$ decides
$q$: it excludes $q$ unless it is negated.  Then
\[
  \Match(\mathcal R,p,t)\;=\;\exists\, q\in\Anc(p):\ q\text{ is excluded by }\mathcal R.
\]
Because an excluded ancestor already makes the whole path excluded, a later
negation cannot re-include something inside it, as in git.
\end{definition}

\begin{definition}[Match below]
\label{def:matchbelow}
$\MatchBelow(\mathcal R,f,p,t)$ is $\Match$ restricted to the prefixes of $p$
that are strictly longer than $f$.  It is used to apply hide rules while listing
a directory: entering a hidden directory on purpose shows its contents.
\end{definition}

\subsection{Laws for a rule set without negation (Deny)}

\begin{law}[Deny monotonicity]
\label{law:deny-monotone}
If $\mathcal R\subseteq\mathcal R'$ then
$\Match(\mathcal R,p,t)\Rightarrow\Match(\mathcal R',p,t)$.  Adding a deny rule
never makes anything accessible.
\Tests{rules/laws_test.go::TestLawDenyIsMonotoneAndOrderIndependent}
\end{law}

\begin{law}[Order independence]
\label{law:deny-order}
If $\mathcal R'$ is a permutation of $\mathcal R$, then
$\Match(\mathcal R,p,t)=\Match(\mathcal R',p,t)$.
\Tests{rules/laws_test.go::TestLawDenyIsMonotoneAndOrderIndependent}
\end{law}

\begin{law}[No negation in deny]
\label{law:no-negation}
Parsing a deny file that contains a negated rule fails.  Together with
Laws~\ref{law:deny-monotone} and~\ref{law:deny-order} this gives the security
reading of a deny file: it is a plain \emph{set} of exclusions.
\Tests{rules/laws_test.go::TestLawDenyRefusesNegation}
\end{law}

\subsection{Laws for every rule set}

\begin{law}[Ancestor closure]
\label{law:ancestor}
If $\Match(\mathcal R,p,\Dir)$ then $\Match(\mathcal R,p\cdot n,t)$ for every
name $n$ and type $t$.  This holds also when $\mathcal R$ contains negated
rules.
\Tests{rules/laws_test.go::TestLawAncestorClosure}
\end{law}

\begin{law}[Fold invariance]
\label{law:fold-invariance}
If $\fold{p}=\fold{p'}$ then $\Match(\mathcal R,p,t)=\Match(\mathcal R,p',t)$:
a path is excluded exactly when any spelling that the file system treats as the
same name is.  In particular no spelling of a denied name (upper case, decomposed
accents, long s, Kelvin sign, sharp s, ligatures) escapes a rule.
\Tests{rules/laws_test.go::TestLawMatchDependsOnlyOnTheFoldedPath; guard/foldalias_test.go}
\end{law}

\begin{law}[Narrowing]
\label{law:narrowing}
$\MatchBelow(\mathcal R,f,p,t)\Rightarrow\Match(\mathcal R,p,t)$.  Ignoring
ancestors can only remove matches.
\Tests{rules/laws_test.go::TestLawMatchBelowNeverAddsMatches}
\end{law}

\begin{law}[Single component]
\label{law:component}
If a pattern contains neither \code{/} nor \code{**}, then $\Hit(r,q,t)$ depends
only on the final component of $q$.  Glob classes never match \code{/}: a
negated class excludes it, and a \code{/} inside a class is a pattern error
(fail closed).
\Tests{rules/laws_test.go::TestLawSlashlessPatternsStayWithinOneComponent; TestSlashInsideAClassIsAParseError}
\end{law}

\begin{law}[Credential folders match wherever they are]
\label{law:credential-anywhere}
The core rules for credential folders (\code{.ssh}, \code{.aws}, \code{.gnupg},
\code{.kube}, \code{.config/gh}, Keychains, and the browser profiles) and files
(\code{.netrc}) are not anchored to the current home directory: a path that
contains one of them at any depth is excluded.  The same secrets exist in a
backup or clone of the home folder, in another account's home, and under a wrong
\code{\$HOME}, where an anchored rule protects nothing.
\Tests{guard/differential\_test.go::TestCredentialFoldersAreDeniedWhereverTheyAre}
\end{law}

\begin{law}[A rule file works as written or fails loudly]
\label{law:rule-file}
For every rule file, either parsing fails with a message, or each rule line means
what a reader would take it to mean.  In particular no accepted line is a rule that
can never match: a leading byte order mark and any line ending are normalized, and
an inline comment, \code{\~{}user/}, a quoted pattern, an invisible character or a
lookalike of the slash are refused.
\Tests{rules/parse\_audit\_test.go}
\end{law}

\begin{law}[The quick test never changes an answer]
\label{law:fast-path}
Matching first asks whether any rule could match, using one combined expression
per kind of rule, and tests slash-free patterns against the last path component
only.  These are optimizations: for every rule set without negation and every
path, the result is that of consulting the rules one by one against the whole
path.
\Tests{rules/laws\_test.go::TestLawFastPathEqualsTheRulesOneByOne}
\end{law}

\begin{law}[Trailing whitespace is not significant]
\label{law:tail}
For every rule $r$ and path $q$, $\Hit(r,q\cdot\!\!\text{``}n\text{''},t)=
\Hit(r,q\cdot\!\!\text{``}n\ \text{''},t)$ for names whose trailing characters
are white space (space, tab, newline).  macOS permits such names and they must
not slip past a rule for \code{*.pem}.
\Tests{rules/rules_test.go; guard/review_test.go}
\end{law}

\subsection{Boundaries}

\begin{boundary}[Hard links]
Rules name paths, so a hard link to a denied file is invisible to them.  The guard
therefore judges a regular file with more than one name by its \emph{identity}
(device and inode): if any name of the file is denied, every name is refused, in
\code{Open} and in listings alike (Law~\ref{law:linked}).  The other names come
from an index of all multi-name files under the roots and under the protected
credential locations, built in the background because visiting every file takes
seconds.  A multi-name file is served only if the index accounts for all of its
names now: it found as many distinct names (by folder identity and folded name) as
the file has links, none is denied, and each is still, when the file is used, a real
path to that file.  Otherwise it is refused: while the index is not ready, for a
name outside every root and protected location or in a folder that could not be
read, after a walk stopped at its limit, and after the names changed.  Within one
listing or search a verdict is reused for at most one second, so a name or size
shown may lag a rename by that long; \code{Open} never reuses it.  Not detected:
a denied name outside every root and protected location, whose file is refused
rather than served.
\end{boundary}

\begin{boundary}[Case classes in ranges]
A positive class range such as \code{[+-0]} can include \code{/} by range even
though the class text does not.  No shipped rule uses such a range; a future
tightening would expand ranges before compiling.
\end{boundary}

\section{The Access Decision}
\label{sec:access}

\subsection{One question, asked once}

Every endpoint asks the same question of the guard: may this path be read?  The
answer is defined here and every operation is then compared with it.

\begin{definition}[Access]
\label{def:ok}
Let $c=\Canon(\Clean(p))$ be the request path made absolute, cleaned, and
reduced to one alias spelling, and let $\rho=\Real(c)$ be the real location of
what is actually opened.  The type $t$ is the type of the opened object.  Then
\[
  \Ok(p)\;\iff\;
  \InRoots(\rho)\ \land\ \neg\Match(\Deny,\rho,t)\ \land\ \neg\Match(\Deny,c,t)
  \ \land\ \Elig(o),
\]
where $o$ is the opened object and $\Elig(o)$ is the hard-link condition: it holds
unless $o$ is a regular file with $n>1$ names, in which case it holds iff the
index accounts for all of them \emph{now}.  That means that the index is ready,
that it found $n$ distinct names (by folder identity and folded name), that no
deny rule matches the real path of any of them, and that each is still a real
path to $o$ (same device and inode, no symbolic link along it).  If the index is
not ready, or any part cannot be established, $\Elig(o)$ is false: unknown means
refused.  Identity is (device, inode); no content is read.
\end{definition}

Without $\Elig$ the formula would be a path-and-root predicate only.  An allowed
\code{ordinary.txt} that has a second name \code{.env} satisfies the first three
conjuncts and is nevertheless refused, which is the point of the fourth
(Law~\ref{law:linked}).

The check is made on the \emph{real} location (so a symbolic link cannot lead
into a denied place), and also on the \emph{spelled} location (so a link
\emph{named} like a denied file is refused even though its target is fine).  It
is therefore conservative: a path is refused if either name is denied.

\begin{definition}[Open, list, search]
$\Open(p)$ returns the opened object, or $\NotFound$, exactly when $\Ok(p)$
fails or the path does not exist, and otherwise a permission or type error that
mentions nothing about denied paths.  $\List(d)$ is defined when $\Open(d)$ is;
it contains an entry $e$ of $d$ exactly when $\Visible(d,e)$:
\begin{align*}
  \Visible(d,e)\iff{}& \neg\Match(\Deny,\rho_d\cdot e,t)\ \land\
      \neg\MatchBelow(\Hide,\rho_d,\rho_d\cdot e,t)\\
  &\land\ (e\text{ a symbolic link}\Rightarrow\Ok\text{-target}(e))\ \land\
      \Elig_\tau(e),
\end{align*}
where $\Elig_\tau$ is $\Elig$ with one difference: within a single listing or
search a verdict about a file may be reused for at most $\tau=1$ second, after
which the names are checked again.  A listing is therefore a snapshot whose
hard-link verdicts can be up to $\tau$ old; $\Open$ always evaluates $\Elig$
afresh,
and $\rho_d=\Real(d)$: children are judged by where the directory really is,
not by the link it was reached through.  $\Search(d,q)$ walks breadth-first from
$d$, entering only visible unlinked directories whose real location is still the
name they were queued under, and reports the visible entries whose folded name
contains $\fold{q}$.  It also reports whether it was \emph{complete}: not if a
folder that was allowed could not be opened or read to the end, or was replaced by a
link before it was entered; denied, vanished and cloud-only folders do not count.
\end{definition}

\subsection{Laws}

\begin{law}[Denied is absent]
\label{law:absent}
If $\neg\Ok(p)$ because of a deny rule, then $p$ is not opened, not listed by
$\List(\Parent(p))$, and not reported by any $\Search$.
\Tests{guard/laws_test.go::TestLawDeniedIsAbsentEverywhere; guard/fuzz_test.go (independent inode and marker oracle)}
\end{law}

\begin{law}[Every spelling the OS resolves is refused]
\label{law:differential}
For every secret $s$ and every spelling $p$ such that the operating system
resolves $p$ to the same object as $s$ (case, folding aliases, decomposed accents,
doubled separators, dot segments, trailing separators, the data-volume and
symbolic-link aliases, named forks), every operation refuses $p$.  The premise is
decided by the OS (same inode), not by \texttt{fsb}, so the law does not depend on
\texttt{fsb}'s own idea of a name.  Dually, every such spelling of an ordinary file
inside the root is served.
\Tests{guard/differential\_test.go (both directions; fails 1,536 times if the folding fix is reverted)}
\end{law}

\begin{law}[Every name of a denied file is denied]
\label{law:linked}
If a regular file has several names and one of them is matched by a deny rule,
then $\Open$ refuses the file under every name, and no listing shows it.  If the
names cannot be determined, or the index does not account for every one of them
(as many distinct names as the file has links, each still a real path to it), the
file is refused.  Identity is the pair
(device, inode), never a hash of the contents, so an ordinary copy of a secret is
not treated as the secret and no file is read to decide.
\Tests{guard/hardlink\_test.go (all four rule kinds, links made after the index was built, links to files outside the roots, refusal while the index builds, bounded walk, names renamed or moved, names outside every location, the retry walk and cooldown, the one-second verdict reuse in listings and searches); e2e ``a hard link to a denied file is refused under any name''}
\end{law}

\begin{law}[Indistinguishable from missing]
\label{law:indist}
For every endpoint $f$, a denied path $p$ and a missing path $q$ receive the
same status and body: $f(p)=f(q)$.  This includes unreadable files inside a
denied folder, which must not answer 403.
\Tests{guard/laws_test.go::TestLawDeniedIsAbsentEverywhere; server/m2_test.go; e2e ``answers 404, with no content, for denied paths on every endpoint''}
\end{law}

\begin{law}[Listed implies openable]
\label{law:listed-open}
If $e$ is listed in $\List(d)$ and is not a dangling link, then $\Open(d\cdot e)$
succeeds, \emph{provided the file system does not change between the two
operations and the operating system's own access and type conditions hold}
(permissions, a regular file or folder).  This is a snapshot property: within
$\tau$ of a change, a listing may still show a multi-name file that a fresh
$\Open$ refuses (the \emph{Hard links} boundary).  The converse is false by
design: hidden entries are not listed but can be opened.
\Tests{guard/laws_test.go::TestLawListedImpliesOpenable}
\end{law}

\begin{law}[Found implies listed]
\label{law:found-listed}
If $\Search$ reports $m$, then $m$ is an entry of $\List(\Parent(m))$, for a
file system that does not change between the two observations (and, for
multi-name files, to within $\tau$).  Search therefore cannot show what browsing
would not, because both go through one function.
\Tests{guard/laws_test.go::TestLawSearchResultsAreListed}
\end{law}

\begin{law}[Hidden is reachable]
\label{law:hidden-reachable}
If $\Match(\Hide,p,t)$ but $\Ok(p)$ holds, then $\Open(p)$ succeeds.  Hide rules
are cosmetic and never restrict direct access.
\Tests{guard/guard_test.go (hidden entries); e2e ``lists the home folder without denied or ignored entries''}
\end{law}

\begin{law}[Containment]
\label{law:containment}
$\Ok(p)\Rightarrow\exists r\in\Roots:\ \Real(p)\text{ is }r\text{ or lies beneath }r$,
where ``is $r$ or lies beneath $r$'' is the identity relation $\InRoots$ of
Law~\ref{law:root-identity}, not a comparison of spellings.  (Folding belongs to
the name and rule model, where over-matching is safe; it is not used for this
\emph{allow} decision.)
\Tests{guard/guard_test.go; guard/firmlink_darwin_test.go}
\end{law}

\begin{law}[Containment is by identity]
\label{law:root-identity}
$\InRoots(\rho)$ holds iff the leading components of $\rho$ that correspond to a
root name \emph{the same directory} as that root (same device and inode).
Comparing spellings would be right on a case-insensitive volume, but on a
case-sensitive one it would admit a sibling folder that differs from the root
only in case.  This is an \emph{allow} decision, so unlike the deny rules it must
not over-match.
\Tests{guard/casesens\_test.go (needs FSB\_CASE\_SENSITIVE\_VOLUME); guard/guard\_test.go::TestInRoots}
\end{law}

\begin{law}[Alias reduction is by name equality]
\label{law:alias-fold}
$\Canon$ recognizes the data-volume prefix and the firmlink directories
component by component, comparing components with $\Fold$, so no spelling of
\code{/System/Volumes/Data/Users} escapes reduction.  Slicing by byte length
does not have this property, because a folded spelling can be longer or shorter.
\Tests{guard/foldalias\_test.go::TestRootOfSlashHonoursDenyRulesThroughFoldedFirmlinkSpellings; guard/firmlink\_darwin\_test.go}
\end{law}

\begin{law}[The session is scoped to a random path]
\label{law:cookie-path}
Browsers scope a cookie by host and never by port, so a cookie for
\code{127.0.0.1} would accompany every request to every other web server on that
address.  All URLs of \texttt{fsb} therefore live under a random 128-bit prefix
$/\pi/$ fixed at launch; the cookie is set with \code{Path=/$\pi$/}; and a
request whose path is not under the prefix is refused exactly like an
unauthenticated one, whatever cookie it carries.  A cookie that is sent only for
paths under $/\pi/$ is never sent to another server, which has no such path.
The frontend forms every request relative to the page's own directory.
\Tests{httpguard/httpguard\_test.go::TestOnlyURLsUnderThePrefixExist, TestSessionCookieIsScopedToThePrefix; e2e ``the session cookie is scoped to fsb's own prefix'' (fails if the path is widened); web/src/lib/laws.test.ts::apiBase}
\end{law}

\begin{law}[Canonical form is stable]
\label{law:canon}
$\Canon(\Canon(p))=\Canon(p)$.
\Tests{guard/laws_test.go::TestLawCanonPathIsIdempotent}
\end{law}

\begin{law}[Bounded work]
\label{law:bounded}
$\Search$ examines at most $V$ directory entries (excluded ones included) and
reports at most $M$ results,
and reports truncation when a bound stopped it; it never follows symbolic links
to directories, so it terminates on cyclic links.
\Tests{guard/search_test.go; e2e (100,000-file directory, optional)}
\end{law}

\subsection{Boundaries}

\begin{boundary}[Time of check]
The file is opened first and its real path is then resolved and confirmed to
name the same object (\code{SameFile}); the decision is taken on the object that
is actually served.  A rename between opening and resolving is treated as a
race and refused.
\end{boundary}

\begin{boundary}[Response timing]
A denied path can be answered slightly faster than a missing one.  Status and
body are identical; timing is not equalized.
\end{boundary}

\begin{boundary}[Other processes]
Anything that runs as the user can read the same files.  \texttt{fsb} is not a
sandbox against local malware.
\end{boundary}

\section{Content-Typed Endpoints}
\label{sec:content}

Three endpoints look inside a file: inline image preview, inline PDF, and
archive listing.  Each is defined by a predicate on the file's \emph{bytes},
never its name, and each factors through $\Open$, so Laws~\ref{law:absent}
and~\ref{law:indist} apply to them unchanged.  A fourth, the Quick Look picture,
hands bytes to another program; Law~\ref{law:quicklook} states what it may be
given.

\begin{definition}[Sniff]
$\Sniff_{\mathrm{img}}$ accepts a byte sequence that begins with the PNG, JPEG,
GIF, or WebP signature; $\Sniff_{\mathrm{pdf}}$ accepts one that begins with
\code{\%PDF-}; $\Sniff_{\mathrm{arc}}$ accepts the start of a zip, a tar
(\code{ustar} at offset 257), or a gzip stream that contains a tar.
\end{definition}

\begin{law}[Content decides]
\label{law:sniff}
The endpoint succeeds only if the corresponding sniff accepts the file's leading
bytes; the file name is never consulted.  A file named \code{a.png},
\code{a.pdf} or \code{a.zip} with other contents is refused with 415, and a file
with a different name but the right contents is accepted.  SVG and HTML are never
accepted.
\Tests{guard/m2_test.go (OpenImage); guard/pdf_test.go; guard/archive_test.go; server/pdf_test.go}
\end{law}

\begin{law}[Cloud-only files are never read]
\label{law:dataless}
For a file whose contents are stored only in the cloud, no endpoint reads it
implicitly: the guard reports it as not downloaded before any read.
\Tests{guard/dataless tests; e2e (hover peek)}
\end{law}

\begin{law}[Archive listing discloses names only]
\label{law:archive-names}
The archive listing contains, for each member, its name, size, type and date,
and never any member data; names are displayed as text with control characters,
Unicode formatting characters and invalid UTF-8 replaced, are cut to a bounded
prefix (4{,}096 bytes, then an ellipsis), and are never used as paths.
\Tests{guard/archive_test.go::TestListArchiveNamesAreDefanged; server/archive_test.go}
\end{law}

\begin{law}[Archive listing is bounded]
\label{law:archive-bounded}
For any file, listing examines at most $E$ tar entries or zip directory records,
streams at most $D$ decompressed bytes of a compressed tar, reads at most $B$
bytes of a zip directory, returns at most $S$ entries, and reports (as
\emph{incomplete}) when a scan limit stopped it.  A \emph{valid} archive that
exceeds a scan limit is listed up to the limit and marked incomplete; it is not
refused.  A zip whose directory was read in full is refused if its record count
disagrees with the archive's end record, which can be checked only then.  The
byte bound $B$ is checked between records, so the last record may cross it by at
most one record's length, and the work stays bounded.  A plain tar is read by
skipping each member's body, so its cost does not grow with member sizes.
\Tests{guard/archive\_test.go::TestListArchiveIsBounded, TestListArchiveStopsAtTheDecompressionLimit, TestZipEntryLimitHoldsForValidArchives, TestZipDirectoryByteBudget, TestZipWithAMisleadingCountIsRefused, TestTarGzLimitAtAHeaderBoundaryIsIncomplete, TestPlainTarListingStillSeeksPastMemberBodies}
\end{law}

\begin{law}[Inline PDF is framed only by the app]
\label{law:pdf-frame}
The PDF response carries \code{X-Frame-Options: SAMEORIGIN} and
\code{frame-ancestors 'self'}, the app shell carries \code{DENY} and
\code{frame-ancestors 'none'}, and the Markdown frame document carries its own
policy that admits exactly its one script and one style by hash.
\Tests{server/pdf_test.go; server/mdframe_test.go; web/mdframe_test.go; e2e PDF and Markdown tests}
\end{law}

\begin{law}[Quick Look sees only what $\Open$ serves]
\label{law:quicklook}
Let $C(p)$ be what Quick Look is given for a request on $p$.  $C(p)$ is defined
only if $\Open(p)$ succeeds, the real path has an iWork or Office extension, and
the file is not cloud-only.  For a file, $C(p)$ is a new file in a private folder
whose bytes are those read from the descriptor $\Open(p)$ returned.  For a
package, $C(p)$ is the set of pairs $(r, b)$ such that $\Open(p/r)$ succeeds on a
regular file whose real path lies under the real path of $p$, and $b$ is the bytes
read from that descriptor.  In both cases $C(p)$ carries no extended attributes and
no resource fork, and it is never the path $p$ itself.  Hence every byte Quick Look
reads is a byte some request could already read through $\Open$, and
Law~\ref{law:absent} extends to it: no denied file and nothing outside the roots
reaches it, whatever $p$ names (an alias, a link, a package with hidden parts).
\Tests{server/quicklook_test.go (private copy, no extended attributes, package parts, a real Finder alias); server/fuzz_test.go::FuzzPathParameter (stand-in Quick Look echoes its input)}
\end{law}

\begin{boundary}[What Quick Look does with the copy]
The law bounds Quick Look's \emph{input}, not its behavior: a flaw in Apple's
generators, or a document that refers to something outside itself, is not
modeled.  The independent review found no generator that loads local files named
inside a document.
\end{boundary}

\begin{boundary}[PDF sandboxing]
A CSP \code{sandbox} cannot be applied to a PDF, because browsers then refuse to
display it.  The protection is instead the byte check, the framing headers, and
the browser's own isolated PDF viewer.
\end{boundary}

\section{Frontend Functions}
\label{sec:frontend}

These functions are pure, so their laws are checked exhaustively enough by
seeded random testing.

\subsection{Row order}

Let $\Cmp_{k,a,g}$ be the row comparison for sort key $k\in\{\text{name, size,
time, kind}\}$, direction $a$, and grouping flag $g$ (``Folders first'').  A
folder has no size and sorts as smaller than any file; a modification time that
cannot be read sorts as the oldest; ties fall back to the name, then to the
exact name, then to the raw characters (a collator that compares digits by value
calls \code{file1} and \code{file01} equal, and the order must not depend on the
input).

\begin{law}[Total preorder]
\label{law:total}
For every $k,a,g$, $\Cmp$ is antisymmetric
($\Cmp(x,y)=-\Cmp(y,x)$), transitive, and total.  This holds even for rows whose
time is unparseable.
\Tests{web/src/lib/laws.test.ts: compareRows is antisymmetric and transitive}
\end{law}

\begin{law}[Input order does not matter]
\label{law:input-order}
For distinct names, $\Sort(\pi(xs))=\Sort(xs)$ for every permutation $\pi$.
(Two names that differ only by normalization cannot coexist in one directory on
APFS, so the premise excludes them.)
\Tests{web/src/lib/laws.test.ts: sortRows is independent of input order}
\end{law}

\begin{law}[Grouping]
\label{law:grouping}
If $g$, every folder precedes every file, for every key and both directions;
grouping does not flip with the direction.
\Tests{web/src/lib/laws.test.ts: Folders first groups folders}
\end{law}

\begin{law}[Reversal]
\label{law:reversal}
If $\neg g$, the descending order is the reverse of the ascending order.
\Tests{web/src/lib/laws.test.ts: descending is the reverse of ascending}
\end{law}

\subsection{Filter}

\begin{law}[Filter agrees with search]
\label{law:filter}
For case-insensitive matching, $\Filter(n,q)\iff \fold{n}\ni\fold{q}$, where the
client's $\Fold$ equals the server's on the aliases of Observation~\ref{obs:aliases}.
The filter box therefore finds the same names that filename search finds.  With
``Match case'' both sides compare after normalization only.
\Tests{web/src/lib/laws.test.ts: the filter agrees with search; guard/matchcase_test.go}
\end{law}

\subsection{Navigation}

\begin{law}[Hash round trip]
\label{law:hash}
$\FromHash(\ToHash(p,s))=(p,s)$ for every path and selection made of names,
including names containing \code{\%}, \code{\#}, \code{?}, \code{\&}, \code{=},
quotes, spaces, and non-ASCII characters.  A literal \code{?} in a path is always
encoded, so the first \code{?} of a hash starts the options.
\Tests{web/src/lib/laws.test.ts: parseHash inverts pathToHash; format.test.ts}
\end{law}

\begin{law}[Go to path is clean and idempotent]
\label{law:goto}
If $\Goto(s)=p$ then $p$ is absolute, contains no \code{.}, \code{..} or empty
component (except the root itself), never climbs above the root, and
$\Goto(p)=p$.
\Tests{web/src/lib/laws.test.ts: resolveGoto is clean and idempotent; goto.test.ts}
\end{law}

\begin{law}[Roots are prefix closed]
\label{law:roots-prefix}
$p$ is within a root $r$ iff $p=r$ or $p$ begins with $r$ followed by \code{/},
so a sibling sharing a prefix (\code{/Users/mex} against \code{/Users/me}) is
outside.  The client check is advisory; the server's $\InRoots$ decides.
\Tests{web/src/lib/laws.test.ts: withinRoots}
\end{law}

\begin{law}[A name is shown as what it is]
\label{law:display}
No character that reorders or hides text (controls, bidirectional overrides and
isolates, line and paragraph separators, the byte order mark, tag characters) is
sent to the screen inside a file name: each is replaced by a visible marker naming
its code point.  A byte that is not part of valid UTF-8, sent as the escape of
Law~\ref{law:wire}, is replaced by a marker naming its value.  Text that needs none
of them, including right-to-left scripts, is unchanged, and links and downloads use
the real name.
\Tests{web/src/lib/laws.test.ts: displayName; web/src/lib/format.test.ts: displayName shows escaped bytes as a marker; e2e ``names that would disguise themselves are shown with visible markers''}
\end{law}

\begin{law}[Names cross the wire whole]
\label{law:wire}
Let $\mathsf{wire}:\Name\to\mathsf{String}$ be the map the server applies to every
name and path it sends: it keeps each maximal run of valid UTF-8 and replaces each
other byte $b$ by the two-part escape $\mathrm{NUL}\cdot\mathsf{hex}(b)$, two uppercase
hexadecimal digits.  Let $\mathsf{q}$ be the client's encoding of a string into a
query value, which turns each escape into the byte $b$ (\code{\%}$\mathsf{hex}(b)$) and
percent-encodes the rest as UTF-8, and let $\mathsf{dec}$ be the server's one decoding
of a query value into bytes.  Then, for all names $n, n'$:
\begin{enumerate}
  \item $\mathsf{wire}(n)=n$ if $n$ is valid UTF-8;
  \item $\mathsf{wire}(n)$ is valid UTF-8;
  \item $\mathsf{wire}(n)=\mathsf{wire}(n')\implies n=n'$ (no name contains NUL, so an
        escape is never part of a name);
  \item $\mathsf{dec}(\mathsf{q}(\mathsf{wire}(n)))=n$, and the same for a path, component
        by component: what the client sends back names the very bytes it was given.
\end{enumerate}
No escape ever reaches the guard, so Laws~\ref{law:deny-monotone} to
\ref{law:credential-anywhere} are stated over real names alone.
\Tests{server/wire\_test.go::TestWireName; TestWireNameIsOneToOne; server/wire\_linux\_test.go::TestNonUTF8NameRoundTrips; web/src/lib/format.test.ts: queryPath sends escaped bytes as themselves and everything else as UTF-8}
\end{law}

\subsection{Markdown links and images}

\begin{law}[Links resolve to normalized absolute paths]
\label{law:link-path}
If $\Link(b,h)$ names a file, the path is absolute and contains no \code{.},
\code{..} or empty component; relative targets are resolved against the
directory $b$ of the Markdown file and cannot climb above the root.
\Tests{web/src/lib/laws.test.ts: resolveLink yields normalized absolute file paths; markdown.test.ts}
\end{law}

\begin{law}[Links and images carry no address]
\label{law:link-inert}
The generated markup has no \code{href} and no \code{src} on links or images
before sanitizing; remote images are replaced by a placeholder, raw HTML is
escaped, and only \code{http} and \code{https} links are ever opened, in a new
tab without a referrer.  Only base64 raster \code{data:} URLs fetched through the
guarded preview endpoint are put into images.
\Tests{web/src/lib/markdown.test.ts; e2e Markdown tests}
\end{law}

\section{The Preview State Machine}
\label{sec:state}

\subsection{States and events}

For a selected entry the preview pane is in one of the states
\[
  \mathsf{idle},\ \mathsf{loading},\ \mathsf{folder},\ \mathsf{image},\
  \mathsf{quicklook},\ \mathsf{pdf},\ \mathsf{archive},\ \mathsf{text},\ \mathsf{binary},\
  \mathsf{empty},\ \mathsf{dataless},\ \mathsf{error}.
\]
The name of the entry gives a \emph{hint} $h\in\{\text{quicklook, image, pdf,
archive, none}\}$; it is only a hint because the server decides (from the bytes,
or for Quick Look from what Quick Look can draw).  Selecting
an entry moves to $\mathsf{loading}$, after a short delay that keeps holding an
arrow key from firing a request per row, and then:
\begin{itemize}
  \item with hint quicklook (checked first, because an iWork package lists as a
        folder), and the entry is not a dangling link, the pane shows
        $\mathsf{quicklook}$ and, if the picture fails to load, goes to
        $\mathsf{folder}$ for a package or falls back to the head request for a
        file;
  \item a folder goes to $\mathsf{folder}$; a dangling link to $\mathsf{error}$;
  \item with hint image, the pane shows $\mathsf{image}$ and, if the server
        refuses the bytes (415), falls back to the head request;
  \item with hint pdf, a one-byte range request asks the server; a
        \code{200}/\code{206} of type \code{application/pdf} gives
        $\mathsf{pdf}$, a \code{409} gives $\mathsf{dataless}$, anything else
        falls back to the head request;
  \item with hint archive, the listing request gives $\mathsf{archive}$, a
        \code{409} gives $\mathsf{dataless}$, anything else falls back to the head
        request;
  \item otherwise the head request gives one of $\mathsf{text}$,
        $\mathsf{binary}$, $\mathsf{empty}$, $\mathsf{dataless}$, or, on failure,
        $\mathsf{error}$.
\end{itemize}

\subsection{Laws}

\begin{law}[Termination]
\label{law:terminate}
From $\mathsf{loading}$ every run reaches a terminal state
(folder, image, quicklook, pdf, archive, text, binary, empty, dataless, or error)
after at most two requests: one hinted request and at most one fall back to the
head request, which never falls back again.
\end{law}

\begin{law}[Freshness]
\label{law:fresh}
The content shown always belongs to the entry that is currently selected: every
request started for a selection, including the fall back after an image or a Quick
Look picture fails to load, is cancelled when the selection changes, and a cancelled request never
changes the state.
\end{law}

\begin{law}[Frame only after confirmation]
\label{law:frame-confirmed}
An inline frame is created for a PDF only in state $\mathsf{pdf}$, which is
entered only after the server has answered \code{application/pdf} for that
path; a Markdown frame is created only for a text head whose name is a Markdown
name, and receives only sanitized HTML.
\Tests{e2e: ``shows a real PDF in the built-in viewer frame, and treats a fake .pdf as text''; Markdown tests}
\end{law}

\begin{law}[Nothing is shown for a denied entry]
\label{law:preview-denied}
Every request of the state machine goes through the endpoints of
Section~\ref{sec:access}; for a denied entry each answers as for a missing one
(Law~\ref{law:indist}), so the pane shows only ``Not available''.
\Tests{e2e: never lists secret files; answers 404}
\end{law}

\begin{boundary}[Termination and freshness are argued, not enumerated]
Laws~\ref{law:terminate} and~\ref{law:fresh} are properties of the component's
control flow.  Freshness was violated by one path (Finding~\ref{find:stale}) and
is now guaranteed by keeping a single cancellation handle per selection; there
is no automated test for that race, because it needs a controllable slow
response.  It is exercised indirectly by the end-to-end tests that select files
quickly.
\end{boundary}

\section{Findings}
\label{sec:findings}

Stating the laws and then checking them found six defects (1--6), and a
subsequent audit of every place that compares names or paths, and of the
environment assumptions behind the rules, found ten more (7--16).  A second
independent review, of the Quick Look endpoint added later, found two more (17--18)
against Laws~\ref{law:absent} and~\ref{law:indist}.  Each was
reproduced by a failing test before it was fixed.  Findings 1, 7, 8, 10 to 15, 17
and 18 are security defects; the rest are consistency defects that fail closed
or are cosmetic, but each is a place where two parts of the program disagreed.

\begin{longtable}{@{}p{0.03\linewidth}p{0.28\linewidth}p{0.23\linewidth}p{0.35\linewidth}@{}}
\toprule
& Defect & Law broken & Fix \\
\midrule
\endhead
1 & \textbf{Deny bypass through full case folding.}  On APFS the name
\code{.\ss h} opens \code{.ssh}, but the matcher folded with simple folding, so
a rule for \code{\~{}/.ssh} did not match it; \code{Open}, \code{Head} and the
listing of the parent all served the denied folder.  Also long s, the Kelvin
sign and the ligatures.
& \ref{law:fold-invariance}, \ref{law:absent}
& Rules, roots and search compare after full folding (\code{rules.Fold}); the alias
set was measured on this file system and is a test.  \label{find:fold} \\
2 & \textbf{Listed but unopenable.}  A plain file named like a \code{dir/} deny
rule (a file called \code{cache}) was shown by \code{List} but refused by
\code{Open}, because \code{Open}'s pre-check assumed every path is a directory.
& \ref{law:listed-open}
& The pre-check treats the final component as a file; the precise check after
opening still refuses real directories. \\
3 & \textbf{Glob class matched across a slash.}  \code{[!a]} compiled to
\code{[\^{}a]}, which matches \code{/}, so a pattern with no slash could match
across directory components.
& \ref{law:component}
& Negated classes exclude \code{/}; a \code{/} inside a class is a pattern error. \\
4 & \textbf{Row order not transitive.}  A modified time that could not be parsed
compared equal to everything, so \code{a}$=$\code{b}, \code{a}$=$\code{c} but
\code{b}$<$\code{c}.  The resulting order depended on the input order.
& \ref{law:total}, \ref{law:input-order}
& Unreadable times sort as the oldest; comparisons no longer use subtraction. \\
5 & \textbf{Filter and search disagreed.}  The filter box folded with simple
lower-casing, so \code{strasse} did not find \code{Stra\ss e}, while filename
search did.
& \ref{law:filter}
& The filter folds with lower-, upper-, lower-casing, which matches full folding
on the measured aliases. \\
6 & \textbf{Stale fallback.}  After an image failed to load, the fall-back request
had its own cancellation handle, so a slow answer could replace the preview of a
newer selection.  \label{find:stale}
& \ref{law:fresh}
& One handle per selection is shared by every request started for it. \\
7 & \textbf{Firmlink prefix matched by byte length.}  The data-volume prefix and
the firmlink directories were recognized by slicing the path to the prefix's
length, so a folded spelling (\code{/S}long-s\code{ystem}, or \code{Data} with
the \code{st} ligature) was cut mid-character and not reduced.  Under
\code{--root /} the unreduced path was inside the root and matched no rule.
& \ref{law:alias-fold}, \ref{law:absent}
& Components are compared one at a time with \code{Fold}. \\
8 & \textbf{Root containment by spelling.}  Containment folded names, so on a
case-sensitive volume a root \code{Public} also served the different folder
\code{public}.  Folding is safe for deny rules but not for an allow check.
& \ref{law:root-identity}, \ref{law:containment}
& Containment is decided by identity: the corresponding leading directory must be
the root itself. \\
9 & \textbf{Client and CLI comparisons.}  Breadcrumbs sliced the path by the
root's character count (wrong when folding changes length) and compared roots
case-sensitively; a link or hash naming a file with another case or accent
selected nothing; the core-rule warning compared rule text exactly.
& \ref{law:roots-prefix}
& All use the same name equality; breadcrumbs split by components. \\
10 & \textbf{Credential rules anchored to one home directory.}  The rules for
\code{\~{}/.ssh}, \code{\~{}/.aws} and the others named the current home only, so
the same folders in a backup or clone of the home folder, in another account's
home, or under a wrong \code{\$HOME} were served.
& \ref{law:credential-anywhere}
& The credential rules match at any depth, like the secret-file patterns. \\
11 & \textbf{A path is not a file (hard links).}  Rules name paths, so a hard link
to a credential file placed anywhere else was served under its innocent name; this
had been recorded only as a limitation.
& \ref{law:absent}, \ref{law:linked}
& A file with several names is judged by identity: a denied name denies them all, in \code{Open} and in listings.  First done for the credential locations only. \\
12 & \textbf{The session cookie followed the host, not the port.}  Cookies are
scoped by host, so the cookie was also sent to every other web server on
\code{127.0.0.1} that the same browser visited, which could replay it.  It had
been recorded as a limitation.
& \ref{law:cookie-path}
& Everything is served under a random per-launch prefix and the cookie's path is
that prefix; requests outside it are refused. \\
13 & \textbf{Hard links to \code{.env}, \code{.pem} and \code{.key} files.}  The
identity check of finding 11 covered only the credential folders, so a second name
for a secret file elsewhere (\code{notes.txt} linked to a \code{.env}) was still
served.
& \ref{law:linked}
& The check covers every denied name under the roots and protected locations, using
a background index of multi-name files; multi-name files are refused until the
index is ready. \\
14 & \textbf{Rule files that load and do nothing.}  A byte order mark broke the
first rule; bare carriage returns joined lines; an inline comment, \code{\~{}user/},
quotes, an invisible character or a lookalike slash produced a rule that never
matched.  All loaded without an error.
& \ref{law:rule-file}
& The byte order mark and line endings are normalized; the rest are errors that
say what is wrong, so fsb does not start. \\
15 & \textbf{Secrets under other names.}  Only known folders and the extensions
\code{.pem} and \code{.key} were denied, so a private key copied out of
\code{.ssh} (\code{id\_rsa}), a certificate store or a password database was served.
& \ref{law:credential-anywhere}
& The names that identify such files are core rules at any depth (\code{.pub} files
stay reachable). \\
16 & \textbf{A name that shows as another name.}  A right-to-left override or a
control character in a file name reordered or hid what the listing displayed, so
\code{invoice}\textit{U+202E}\code{txt.exe} read as a text file.
& \ref{law:display}
& Deceptive characters are replaced, for display only, by a visible marker; real
right-to-left text is untouched and downloads keep the real name. \\
17 & \textbf{A path that means more to another program.}  Quick Look was given the
real path of a document that passed every check, but a Finder alias named
\code{report.docx} is a regular file that Quick Look resolves to its target, so it
drew a file in a denied folder, or outside the roots.
& \ref{law:absent}
& Quick Look is given only a private copy of bytes read through $\Open$
(Law~\ref{law:quicklook}); a real alias is the regression test. \\
18 & \textbf{A package answered differently for a hidden part.}  A package document
holding a denied file was refused (404) while a clean one answered 200, although
listings show both alike.
& \ref{law:indist}
& Parts a listing leaves out are left out of the copy, and the package answers as a
clean one. \\
\bottomrule
\end{longtable}

\subsection{Observations that are not defects}

\begin{itemize}
  \item A dangling symbolic link is listed (marked as broken) but cannot be
        opened.  Law~\ref{law:listed-open} excludes it.
  \item Hidden entries are not listed but can be opened, and a hidden directory
        that is entered on purpose shows its contents.
        (Law~\ref{law:hidden-reachable}.)
  \item A symbolic link \emph{named} like a denied file is refused even when its
        target is harmless.  Definition~\ref{def:ok} checks both names.
  \item The macOS by-inode path \code{/.vol/<device>/<inode>} was probed as a
        possible alias that would evade path rules; on the tested system the
        directory is empty and the path cannot be opened.
\end{itemize}

\subsection{Defects found afterwards by review}

Reviews of the code (October 2026) found defects that these laws did not, each in
a place the laws state only loosely: a hard-link verdict served for a file whose
names the index had not seen, or whose names changed (Law~\ref{law:linked} now says
what ``determined'' means); a search that entered a folder swapped for a link after it
was queued, and that looked complete when a folder could not be opened (the search
definition now says so); a collator that left two names in input order (the row order
now ends on the raw characters); and several limits that did not bind (a count that could
never trip, a work limit that ignored excluded entries).  The design specification's
history lists them; each has a test that fails without its fix.  Later reviews of
opening the browser on Linux found that the launch URL, in the opener's command line,
let another user take the session; Law~\ref{law:launch-owner} was then stated, and its
review found that a used token still caused a lookup and a warning on every replay (the
law's first case now comes before the owner), and that the tests could not tell the
client's end of a connection from \texttt{fsb}'s own, which would have satisfied the
law's form while letting anyone through.

\subsection{What the method did and did not find}

Findings 1 to 3 come from the rule and access laws and are about names: two
spellings, two types, two segments treated as one.  Findings 4 to 6 come from
the frontend laws.  Findings 7 to 9 came from asking of every comparison in the
code whether it assumes that the same name means the same bytes, and from
checking each answer against the real file system (including a case-sensitive
volume).  Findings 1, 7 and 8 share one lesson: name equality is a property of
the volume, so deny checks may over-match but allow checks must decide by
identity.  Finding 17 adds a second: a path that the guard has judged is safe
only for code that opens it the way the guard did; another program may see more in
it, so what crosses that boundary must be bytes, not a name.  The laws that held on first check (monotonicity, order
independence, ancestor closure, narrowing, search implies listed, denied is
absent, the round trips) are as informative as the failures, since they turn
long-standing assumptions into checked statements.

\section{Refinement and Summary}
\label{sec:refinement}

\subsection{From the model to the code}

\begin{center}
\small
\begin{tabular}{@{}p{0.30\linewidth}p{0.62\linewidth}@{}}
\toprule
Abstract & Concrete \\
\midrule
$\Fold$ & \code{rules.Fold} (Go); \code{fold} in \code{web/src/lib/filter.ts} \\
$\Match$, $\MatchBelow$ & \code{Set.Match}, \code{Set.MatchBelow} in \code{internal/rules} \\
$\Deny$, $\Hide$ & the two \code{rules.Set} values held by \code{guard.Guard} \\
$\Canon$, $\Real$ & \code{canonPath}, \code{filepath.EvalSymlinks} in \code{internal/guard} \\
$\Ok$, $\Open$ & \code{Guard.open}, with \code{violation} and \code{deniedByAncestor} \\
$\List$, $\Visible$ & \code{Guard.ListFunc}, \code{Guard.entry} \\
$\Search$ & \code{Guard.Search} (uses \code{entry}) \\
$\Sniff$ & \code{sniffImage}, \code{OpenPDF}, \code{ListArchive} \\
$\Cmp$, $\Sort$ & \code{compareRows}, \code{sortRows} in \code{sort.ts} \\
$\Filter$ & \code{matchesName} in \code{filter.ts} \\
$\ToHash$, $\FromHash$ & \code{pathToHash}, \code{parseHash} in \code{format.ts} \\
$\Goto$ & \code{resolveGoto} in \code{goto.ts} \\
$\Link$ & \code{resolveLink} in \code{markdown.ts} \\
\bottomrule
\end{tabular}
\end{center}

Every HTTP-facing endpoint reaches the file system only through \code{Guard}; a
test fails the build if an HTTP-facing package imports \code{os}, \code{io/fs},
\code{syscall} or \code{path/filepath}.  This is what allows
Laws~\ref{law:absent} and~\ref{law:indist} to be stated once for the guard rather
than once per endpoint.

\subsection{How to run the checks}

\begin{lstlisting}[language=bash]
go test -race ./...                 # rule, guard and server laws
cd web && npm test                  # frontend laws
node --test e2e/chrome.test.mjs     # end-to-end (Safari: e2e/safari.test.mjs)
make -C specs/algebra                     # build this document
\end{lstlisting}

\subsection{Summary}

\begin{itemize}
  \item The deny rules form a set of exclusions: monotone, order independent,
        closed under descent, and invariant under every spelling that the file
        system equates.
  \item Every endpoint agrees with one access decision; listing implies
        openable, search implies listed, and denied paths are absent and look
        missing.
  \item Content-typed endpoints decide from bytes, disclose names only, and are
        bounded.
  \item The frontend order is a total preorder independent of input order, and
        the filter, hash and link functions are the exact counterparts of their
        server-side definitions.
  \item Sixteen defects were found and fixed; the laws that hold are recorded with
        the checks that establish them.
\end{itemize}

The document is intended to be read alongside \code{SECURITY.md} and the product
functional and design specifications, which sit in the enclosing \code{specs} folder: those state what is promised and how it is built, and this states, precisely, the
properties from which the promise follows.


\end{document}
